\chapter{Algorithmes et programmation}

\section{Variable, affectation et types de données}

\begin{defin}[Variable et affectation]
	En programmation, une \textbf{variable} est un espace mémoire nommé servant à stocker une donnée.

	L'\textbf{affectation} est l'opération qui consiste à attribuer une valeur à une variable.

	En Python, elle est réalisée à l'aide du symbole \lstinline|=|.
\end{defin}

\begin{remarque}
	Dans une instruction d'affectation ~\lstinline|x = expression|~, l'évaluation de l'expression située à droite du ~\lstinline|=|~ est effectuée en premier, puis le résultat est stocké dans la variable ~\lstinline|x|.
\end{remarque}

\begin{ex}
	\nline
	\begin{lstlisting}
a=5
b=a+2
a=a+1
\end{lstlisting}

	\begin{itemize}
		\item \lstinline|a=5| : affecte la valeur entière $5$ à la variable \lstinline|a|.
		\item \lstinline|b=a+2| : évalue $5 + 2$, puis affecte la valeur $7$ à la variable \lstinline|b|.
		\item \lstinline|a=a+1| : augmente de $1$ la valeur contenue dans \lstinline|a| (opération d'incrémentation).
	\end{itemize}
\end{ex}

\begin{prop}[Principaux types de données en Python]
	En Python, chaque variable possède un type déterminé par la valeur qu'elle contient :
	\begin{itemize}
		\item \textbf{Entiers} (\lstinline|int|) : nombres entiers relatifs (ex. : \lstinline|3|, \lstinline|-12|).
		\item \textbf{Flottants} (\lstinline|float|) : nombres décimaux (ex. : \lstinline|3.14|, \lstinline|-0.5|).
		\item \textbf{Chaînes de caractères} (\lstinline|str|) : suites de caractères délimitées par des guillemets (ex. : \lstinline|"Maths"|, \lstinline|'Python'|).
		\item \textbf{Booléens} (\lstinline|bool|) : valeurs de vérité, soit \lstinline|True| (Vrai), soit \lstinline|False| (Faux).
	\end{itemize}
\end{prop}

\begin{attention}
	Le séparateur décimal en Python est le point \lstinline|.| et non la virgule \lstinline|,|.

	De plus, la chaîne de caractères \lstinline|"5"| de type \lstinline|str| est différente du nombre entier \lstinline|5| de type \lstinline|int|.
\end{attention}

\begin{prop}[Opérateurs arithmétiques]
	Soient \lstinline|a| et \lstinline|b| deux variables numériques. Les principaux opérateurs sont :
	\begin{multicols}{2}
		\begin{itemize}
			\item \lstinline|a + b| : Addition
			\item \lstinline|a - b| : Soustraction
			\item \lstinline|a * b| : Multiplication
			\item \lstinline|a / b| : Division exacte (résultat \lstinline|float|)
			\item \lstinline|a // b| : Quotient de la division entière
			\item \lstinline|a % b| : Reste de la division entière (modulo)
			\item \lstinline|a ** b| : Puissance ($a^b$)
		\end{itemize}
	\end{multicols}
\end{prop}

\newpage

\begin{methode}[Échanger les valeurs de deux variables]
	Pour échanger les contenus de deux variables \lstinline|a| et \lstinline|b|, il est nécessaire de passer par une variable temporaire \lstinline|temp| pour ne pas perdre de donnée, ou d'utiliser l'affectation simultanée propre à Python.

	\begin{correction}
		\nline
		{
			\setlength{\columnseprule}{0.6pt}
			\begin{multicols}{2}
				\noindent\textbf{Méthode avec variable temporaire :}
				\begin{lstlisting}[language=Python]
temp = a
a = b
b = temp
\end{lstlisting}
				\columnbreak
				\noindent\textbf{Méthode propre à Python :}
				\begin{lstlisting}[language=Python]
a, b = b, a
	
	
\end{lstlisting}
			\end{multicols}
		}
	\end{correction}
\end{methode}

\section{Structure conditionnelle}

\begin{defin}[Structure conditionnelle]
	Une \textbf{structure conditionnelle} permet d'exécuter un bloc d'instructions uniquement si une condition donnée (une expression booléenne) est vérifiée.
\end{defin}

\begin{prop}[Opérateurs de comparaison]
	Les conditions s'expriment à l'aide des opérateurs de comparaison suivants :
	\begin{multicols}{3}
		\begin{itemize}
			\item \lstinline|==| : Égal à
			\item \lstinline|!=| : Différent de
			\item \lstinline|<| : Strictement inférieur
			\item \lstinline|>| : Strictement supérieur
			\item \lstinline|<=| : Inférieur ou égal
			\item \lstinline|>=| : Supérieur ou égal
		\end{itemize}
	\end{multicols}
\end{prop}

\begin{attention}
	Il ne faut pas confondre le symbole d'affectation \lstinline|=| avec le test d'égalité \lstinline|==|.
\end{attention}

\begin{defin}[Syntaxe des blocs \lstinline|if / elif / else|]
	En Python, les blocs d'instructions sont délimités par des deux-points \lstinline|:| et par l'\textbf{indentation} (décalage vers la droite systématique d'une tabulations ou $4$ espaces).

	\begin{lstlisting}[language=Python]
if condition_1:
    # Instructions si condition_1 est vraie
elif condition_2:
    # Instructions si condition_1 est fausse et condition_2 vraie
else:
    # Instructions si aucune condition n'est vraie
\end{lstlisting}
\end{defin}

\begin{methode}[Mise en œuvre d'une structure conditionnelle]
	On souhaite déterminer si un nombre réel $x$ appartient à l'intervalle $\intff{2}{5}$.

	\begin{correction}
		\nline
		\begin{lstlisting}[language=Python]
if x >= 2 and x <= 5:
    print("x appartient a l'intervalle [2 ; 5]")
else:
    print("x n'appartient pas a l'intervalle [2 ; 5]")
\end{lstlisting}
	\end{correction}
\end{methode}

\begin{remarque}[Opérateurs logiques]
	Les conditions complexes peuvent être combinées à l'aide des opérateurs logiques :
	\begin{itemize}
		\item \lstinline|and| (ET) : vrai si toutes les conditions sont vraies simultanément.
		\item \lstinline|or| (OU) : vrai si au moins l'une des conditions est vraie.
		\item \lstinline|not| (NON) : inverse l'état logique d'une condition.
	\end{itemize}
\end{remarque}


\section{Structures itératives (Boucles)}

\begin{defin}[Boucle bornée \texttt{for}]
	Une \textbf{boucle bornée} répète un bloc d'instructions un nombre déterminé à l'avance de fois.

	En Python, elle s'écrit avec la commande \lstinline|for| combinée avec \lstinline|range()|.
\end{defin}

\begin{prop}[Fonction \texttt{range}]
	La fonction \lstinline|range()| génère une séquence d'entiers :
	\begin{itemize}
		\item \lstinline|range(n)| génère les $n$ entiers de $0$ à $n-1$.
		\item \lstinline|range(a, b)| génère les entiers de $a$ à $b-1$.
		\item \lstinline|range(a, b, k)| génère les entiers de $a$ à $b-1$ par pas de $k$.
	\end{itemize}
\end{prop}

\begin{ex}
	Le code suivant calcule la somme $S = 1 + 2 + 3 + \dots + 10$ :
	\begin{lstlisting}[language=Python]
S = 0
for i in range(1, 11):
    S = S + i
print(S)
\end{lstlisting}
\end{ex}

\begin{defin}[Boucle non bornée \texttt{while}]
	Une \textbf{boucle non bornée} répète un bloc d'instructions tant qu'une condition reste vraie.

	Le nombre d'itérations n'est pas connu à l'avance.
	\begin{lstlisting}[language=Python]
while condition:
    # Instructions a repeter
\end{lstlisting}
\end{defin}

\begin{attention}[Risque de boucle infinie]
	Dans une boucle \lstinline|while|, la condition doit obligatoirement devenir fausse à un moment donné.

	Si la variable de contrôle n'est pas modifiée dans le corps de la boucle, le programme boucle indéfiniment.
\end{attention}

\begin{methode}[Utiliser une boucle \lstinline|while|]
	On place la somme de $1000$ € sur un compte à un taux d'intérêts annuel de $5\,\%$.

	On souhaite déterminer au bout de combien d'années le capital dépassera $1500$ €.

	\begin{correction}
		\nline
		\begin{lstlisting}[language=Python]
capital = 1000
annees = 0
while capital <= 1500:
    capital = capital * 1.05
    annees = annees + 1
print("Nombre d'annees :", annees)
\end{lstlisting}
	\end{correction}
\end{methode}

\section{Fonctions et modularité}

\begin{defin}[Fonction]
	En programmation, une \textbf{fonction} est un bloc d'instructions nommé qui effectue une tâche précise.

	Elle peut recevoir des paramètres d'entrée (arguments) et renvoyer un résultat à l'aide du mot-clé \lstinline|return|.
\end{defin}

\begin{defin}[Syntaxe d'une fonction]
	\nline
	\begin{lstlisting}[language=Python]
def nom_de_la_fonction(parametre1, parametre2):
    # Bloc d'instructions
    resultat = ...
    return resultat
\end{lstlisting}
\end{defin}

\begin{methode}[Définir et appeler une fonction mathématique]
	Soit la fonction mathématique $f: x \mapsto 3x^2 - 5x + 2$.

	Créer une fonction Python nommée \lstinline|f| prenant en paramètre un nombre \lstinline|x| et renvoyant son image par $f$.

	\begin{correction}
		\nline
		\begin{lstlisting}[language=Python]
def f(x):
    y = 3 * x**2 - 5 * x + 2
    return y

# Appel de la fonction pour x = 4
image = f(4)
print(image)
\end{lstlisting}
	\end{correction}
\end{methode}

\begin{remarque}
	L'instruction \lstinline|return| interrompt immédiatement l'exécution de la fonction et renvoie la valeur spécifiée. Si une fonction ne contient pas d'instruction \lstinline|return|, elle renvoie la valeur \lstinline|None|.
\end{remarque}

\begin{methode}[Algorithme de calcul de distance dans le plan]
	Créer une fonction \lstinline|distance(xA, yA, xB, yB)| qui renvoie la distance entre deux points $A(x_A, y_A)$ et $B(x_B, y_B)$ dans un repère orthonormé :
	\[
		AB = \sqrt{(x_B - x_A)^2 + (y_B - y_A)^2}
	\]

	\begin{correction}
		\nline
		\begin{lstlisting}[language=Python]
from math import sqrt

def distance(xA, yA, xB, yB):
    d = sqrt((xB - xA)**2 + (yB - yA)**2)
    return d
\end{lstlisting}
	\end{correction}
\end{methode}
